\section{Two Diagnostic Cases}
\label{sec:cases}

The cases in this section are not used as proof of the model. They illustrate how the framework changes the questions an empirical researcher asks.

\subsection{South Korean footwear: decline is not enough to establish supersession}

South Korea's footwear industry, centered heavily in Busan, became a major export cluster and then contracted sharply as wages and land costs rose and labor-intensive production shifted to lower-cost locations. Lim documents the restructuring problem already in the early 1990s. The OECD's territorial review of Busan describes a steep fall in footwear exports and substantial relocation of production abroad, while also noting continued local specialization. Later historical work describes attempts at automation, own-brand development, specialization, and niche production during the industry's decline.

For the present framework, the important point is methodological. Sectoral decline is not automatically evidence of developmental failure, because rising wages and changing comparative advantage can make mass labor-intensive assembly a poor use of accumulated capability. But the opposite inference is also invalid: the fact that Korea later became a high-technology economy does not prove that footwear capabilities migrated directly into semiconductors or automobiles.

The case therefore requires tracing what survived the old configuration. Relevant evidence would include worker transitions, design and engineering routines, supplier capabilities, logistics and export competence, retained brand or niche specialization, capital redeployment, and the destinations of firms that moved production abroad. The Korean case is useful precisely because it may contain all three outcomes at once: destroyed capability, retained and upgraded footwear capability, and recombinable organizational or export competence.

\subsection{FATE and synthetic rubber in Argentina: a visible de-accretion edge}

A sharper negative-linkage example appeared in Argentina in 2026. FATE, a long-standing Argentine tire producer, ceased activity at its Virreyes/San Fernando plant on 18 February 2026. The company's own board communiqu\'e, reproduced contemporaneously by several news outlets, stated that FATE was ceasing industrial activity and terminating its employment contracts after more than eight decades of operation. Contemporary reporting placed the directly affected workforce at roughly 920 employees and described weak demand, competitiveness pressure, rising imports, and labor conflict as part of the surrounding context. The framework does not presume that the plant should therefore have been protected. The relevant question is what happened to the productive network around it.

On 22 July 2026 Pampa Energ\'ia announced the cessation of synthetic-rubber production at its Puerto General San Mart\'in complex. A company statement reproduced in contemporary reporting said changes in the domestic market had made continuation of the business unviable; reporting consistently identified FATE as a principal local customer and connected the demand collapse after FATE's closure to the decision. Roughly 130 workers in the affected operation were reported as subject to redeployment efforts.

This sequence is a concrete candidate example of Proposition~\ref{prop:deaccretion}: a downstream buyer exits; demand for an upstream specialized input falls; the upstream relation can cross below minimum viable scale. The observed chronology and the reported customer relationship support the existence of the edge, but they do not by themselves estimate its full causal magnitude. Nor do they establish the aggregate welfare cost of FATE's closure or prove that all synthetic-rubber capability was permanently destroyed. They establish that the relevant unit of analysis is larger than the exiting factory.

An empirical evaluation should therefore ask at least four separate questions. First, what consumer or downstream efficiency gains came from cheaper or more abundant tire imports? Second, which FATE-specific capabilities were destroyed, retained, or redeployed? Third, which upstream capabilities lost viability because the buyer disappeared? Fourth, did alternative domestic, regional, or export demand emerge quickly enough to preserve those capabilities? The model permits the answer to be mixed: static allocative gains and developmental capability losses can occur simultaneously.

This is why the paper does not equate protection with development or liberalization with development. The diagnostic is whether a policy transition creates a viable path through activation, capability accumulation, graduation, and subsequent structural change without unnecessarily destroying valuable productive state.
